\section{Introduction}
\label{sec:introduction}

A code execution sandbox is now a product you buy by the API call. E2B,
Daytona, Runta, Modal and others expose the same four verbs --- create a
sandbox, run a command in it, suspend it, destroy it --- and the workloads
buying them are AI coding agents and reinforcement-learning training loops.
Both run untrusted, model-generated code, so both need an isolation boundary;
both run it repeatedly against evolving state, so neither fits the stateless
invocation that serverless platforms were built around.

The research community has a mature answer for serverless. SeBS established
what a serverless benchmark suite must do: pick representative applications,
build one framework that deploys and measures them across providers, and fix a
methodology for performance, cost and reproducibility~\cite{sebs2021}.
SeBS-Flow then showed what justifies a \emph{new} suite rather than more
benchmarks in the old one: a new execution model that the existing tests cannot
characterise, in its case function orchestration rather than single
functions~\cite{sebsflow2025}.

This paper argues that cloud code execution sandboxes are such a model, and
measures them.

\subsection*{Why existing benchmarks do not transfer}

Four properties separate a sandbox session from a function invocation, and each
one breaks a measurement that FaaS benchmarks rely on.

\begin{itemize}
  \item \textbf{It is interactive.} A session is a sequence of commands over a
  warm environment, not one call. The cost that matters is the per-command
  overhead on a sandbox that already exists, which a cold-start number does not
  predict.
  \item \textbf{It is stateful.} The filesystem survives between commands, and
  on some providers across a suspend. Correctness is therefore part of the
  measurement: a resume that silently cold-boots is fast and useless.
  \item \textbf{It is reset, not re-invoked.} A reinforcement-learning
  environment returns to a start state thousands of times per training run.
  Depending on the provider and the environment, that reset is an in-process
  operation, a sandbox recreate, or a snapshot restore. Reporting one number
  across the three compares nothing.
  \item \textbf{It is multi-tenant at scale.} A rollout fleet runs thousands of
  sandboxes concurrently. What a trainer waits on is not single-sandbox speed
  but sustainable provisioning throughput and the slowdown caused by
  neighbours.
\end{itemize}

\subsection*{Contributions}

\begin{description}
  \item[C1: A workload model.] We characterise stateful, interactive code
  execution: session length, command duration distribution, the split between
  environment preparation and useful work, reset frequency, state size, and
  arrival pattern (\S\ref{sec:background}). \needsdata{the characterisation
  itself; without it the benchmark design is an engineering guess}
  \item[C2: A benchmark specification.] Five performance dimensions ---
  lifecycle, execution, compute, storage, network --- crossed with two
  execution scales, single and cluster, plus four end-to-end application
  workloads and a reinforcement-learning profile built from the same
  dimensions (\S\ref{sec:design}).
  \item[C3: A provider-agnostic framework.] An open-source Go implementation
  with a deliberately small provider contract, capability reporting that
  distinguishes \emph{unsupported} from \emph{slow}, and a timing method that
  separates a provider's command transport from the sandbox behind it
  (\S\ref{sec:implementation}).
  \item[C4: A cross-provider evaluation.] Measurements across three hosted
  providers answering four research questions (\S\ref{sec:evaluation}).
  \needsdata{all of it}
\end{description}

The framework and every configuration file in this paper are at
\url{https://github.com/jiweiyuan/sandboxbench}.
